% !PRINTED 554
\weberpage{554}
%<<CONT
und $\chi$ ein ganzes Functional. Nun wählen wir (nach \S.\,143, 3.) eine durch $\chi$ theilbare Zahl $\beta$ so, dass $\beta:\chi$ relativ prim zu $\omega$ wird, und setzen

\weq{3}{\psi\chi = \beta.}

Da jetzt $\varphi:\psi$ eine Zahl ist, so ist $\psi$ mit $\varphi$ äquivalent, und $\psi$ ist ein zu $\omega$ theilerfremder Repräsentant der Classe $C^1$).

\S.\,153.

\gesperrt{Die Classenzahl des Körpers $\Omega$.}

Wir kommen nun zum Beweise des wichtigen Satzes:

1. \gesperrt{Die Anzahl der Idealclassen eines Körpers $\Omega$ ist endlich.}

Dieser Satz ist gleichbedeutend mit dem folgenden:

2. \gesperrt{In jeder Classe giebt es ganze Functionale, deren absolute Norm eine bestimmte endliche, nur von der Natur des Körpers $\Omega$ abhängige Zahl nicht übersteigt.}

Denn weil jedes ganze Functional ein Factor seiner absoluten Norm ist, und jede ganze Zahl nur eine endliche Anzahl von nichtassociirten Factoren hat, so giebt es, von den associirten abgesehen, nur eine endliche Zahl von ganzen Functionalen, die eine gegebene Zahl zur absoluten Norm haben. Wenn nun bewiesen werden kann, dass in jeder Classe ein ganzes Functional vorkommt, dessen absolute Norm unter einer endlichen Zahl

\begin{footnotesize}
$^1$) Wir wollen hier im Vorübergehen auf eine Analogie der Aequivalenz der Ideale hinweisen. Die ganze Theorie der algebraischen Zahlen lässt sich mit den nothwendigen Modificationen übertragen auf die Theorie der algebraischen Functionen, die wieder ihren geometrischen Ausdruck in der Theorie der algebraischen Curven findet. Den Idealen entsprechen dann Punktsysteme auf einer festen Grundcurve und den Hauptidealen volle Schnittpunktsysteme der Grundcurve mit einer anderen algebraischen Curve. Wenn sich zwei Punktsysteme zu einem vollen Schnittpunktsystem ergänzen, was in der obigen Formel $\varphi\chi = \alpha$ seinen Ausdruck finden würde, so wird von den beiden Punktsystemen $\varphi$, $\chi$ jedes der \gesperrt{Rest} des anderen genannt. Zwei Punktsysteme, die, wie $\varphi$, $\psi$, denselben Rest haben, werden in der Geometrie corresidual genannt. Dieser Begriff entspricht also der Aequivalenz. (Vergl. \gesperrt{Brill} u. \gesperrt{Nöther}, Mathem. Annalen, Bd. 7. \gesperrt{Salmon}, \glqq higher plane Curves\grqq{}, deutsch von \gesperrt{Fiedler}.)
\end{footnotesize}
%CONT>>
